% Laget automatisk av elevark/lag_elevark.py fra opplegg/video-fall.html. Endringer her blir overskrevet.
\documentclass[11pt]{article}
\usepackage{elevark}
\begin{document}
\elevarkhode{Fysikk 1 · Bevegelse og krefter}{Videoanalyse av fall}{https://fysikklab.no/opplegg/video-fall.html}{Opplegget på nett}

Slipp en ball og film fallet. I Tracker får dere posisjonen i hvert bilde. Er bevegelsen jevnt akselerert, og hvor stor er tyngdeakselerasjonen?

\section*{Hensikt}

\begin{itemize}
\item Bruke videoanalyse til å måle posisjon og tid.
\item Vise at fallet har konstant akselerasjon, og bestemme $g$.
\item Vurdere feilkilder i videoanalyse.
\end{itemize}

\section*{Utstyr}

\begin{itemize}
\item En liten, tung ball, for eksempel en golfball, sprettball eller stålkule.
\item Mobil eller kamera på stativ, helst med høy bildefrekvens (60 bilder per sekund eller sakte film).
\item En meterstokk eller målestokk som henger loddrett ved siden av fallet.
\item Tracker (gratis fra physlets.org/tracker).
\end{itemize}

\section*{Framgangsmåte}

\begin{enumerate}
\item Heng målestokken loddrett. Sett kameraet på stativ et par meter unna, med linsa omtrent midt på fallhøyden og vinkelrett på veggen.
\item Slipp ballen fra ro rett foran målestokken, i samme avstand fra kameraet. Film hele fallet.
\item I Tracker: Sett kalibreringsstaven på målestokken, og plasser aksene med $y$ oppover. Lag en punktmasse og klikk på midten av ballen i hvert bilde, eller bruk autotracker.
\item Kopier kolonnene $t$ og $y$ og lim dem inn i verktøyet på nettsiden.
\item Tilpass en andregradsfunksjon. Akselerasjonen er to ganger tallet foran $t^{2}$.
\end{enumerate}

Uten egen video: Tracker har et eksempel med en ball som kastes, «BallToss», som læreren kan dele. Dere kan også øve med simuleringen Videoanalyse uten video\footnote{\url{fysikklab.no/sim/videoanalyse.html}}.

\section*{Oppgaver}

\begin{oppgaveliste}
\oppgave \fase{gjett}{Gjett}Hvordan tror du grafene for posisjon og fart ser ut? Last inn eksempeldataene og sjekk. Bruk «Vis fart».
\svar
\oppgave \fase{test}{Mål}Tilpass en andregradsfunksjon til hele fallet. Hva blir $g$? Hvorfor er tallet foran $t^{2}$ halvparten av akselerasjonen?
\svar
\oppgave \fase{test}{Undersøk}Prøv å tilpasse en rett linje i stedet. Hvordan ser du på grafen og $R^{2}$ at modellen er dårlig?
\svar
\oppgave \fase{forklar}{Vurder}Eksempeldataene har litt tilfeldig «klikkestøy». Hvor mye endrer $g$ seg hvis du bare bruker første halvdel av fallet? Hvorfor blir resultatet sikrere med hele fallet?
\svar
\oppgave \fase{forklar}{Vurder}Hvilke feilkilder er det i deres egne målinger? Tenk på perspektiv, uskarpe bilder, kalibrering og luftmotstand.
\svar
\oppgave \fase{utfordring}{Utfordring}Gjør forsøket med en ballong eller en lett plastball. Hvordan ser fartsgrafen ut nå? Sammenlign med Fall med luftmotstand\footnote{\url{fysikklab.no/sim/luftmotstand.html}}.
\svar
\end{oppgaveliste}

\end{document}
